\documentclass[11pt,a4paper,oneside]{report}

\usepackage[utf8]{inputenc}
\usepackage{amsmath,amssymb,amsthm}
\usepackage{geometry}
\geometry{a4paper, margin=2.5cm, headheight=28pt}
\usepackage{booktabs}
\usepackage{xcolor}
\usepackage{tcolorbox}
\usepackage{titlesec}
\usepackage{fancyhdr}
\usepackage{array}
\usepackage{hyperref}

\hypersetup{
    colorlinks=true,
    linkcolor=navy,
    citecolor=navy,
    urlcolor=teal,
    pdftitle={Quantum Analysis Framework (QAF): Libro Monográfico Completo},
    pdfauthor={Antoni Bancells, Juan Manuel Goenaga}
}

% Palette
\definecolor{navy}{HTML}{1A365D}
\definecolor{teal}{HTML}{2B6CB0}
\definecolor{darkgray}{HTML}{2D3748}
\definecolor{lightgray}{HTML}{F7FAFC}
\definecolor{accent}{HTML}{DD6B20}

% Header and Footer
\pagestyle{fancy}
\fancyhf{}
\rhead{\textcolor{darkgray}{\small\leftmark}}
\lfoot{\textcolor{darkgray}{\small Quantum Analysis Framework (QAF)}}
\rfoot{\textcolor{darkgray}{\small Página \thepage}}
\renewcommand{\headrulewidth}{0.4pt}
\renewcommand{\footrulewidth}{0.4pt}

% Title Formatting
\titleformat{\chapter}[display]
  {\normalfont\huge\bfseries\color{navy}}{\chaptertitlename\ \thechapter}{20pt}{\Huge}
\titleformat{\section}
  {\normalfont\Large\bfseries\color{teal}}{\thesection}{1em}{}
\titleformat{\subsection}
  {\normalfont\large\bfseries\color{darkgray}}{\thesubsection}{1em}{}

\begin{document}

% Front Matter
\begin{titlepage}
    \centering
    \vspace*{2cm}
    
    {\Huge \bfseries \color{navy} Quantum Analysis Framework\\[0.5em] (QAF)}
    
    \vspace{1.5cm}
    {\Large \itshape Teoría de Decisiones No Conmutativas, Redes Tensoriales e Inteligencia Artificial Metacognitiva a Escala}
    
    \vspace{2.5cm}
    \begin{tcolorbox}[colback=lightgray,colframe=navy,arc=3mm,width=0.85\textwidth,halign=center]
        \Large \bfseries Monografía Técnica y Científica
    \end{tcolorbox}
    
    \vfill
    
    {\large \textbf{Autores:} Antoni Bancells \& Juan Manuel Goenaga}\\[0.5em]
    {\small Colaboración Universitaria: Universidad de los Andes \& UNIR}\\[2cm]
    
    {\small \today}
\end{titlepage}

\tableofcontents
\clearpage

% ==========================================
% PARTE I
% ==========================================
\part{La Crisis de la Decisión Clásica y el Espacio No Conmutativo}

\chapter{La Niebla de la Complejidad y la Necesidad de QAF}
\label{cap:introduccion}

\section{El Colapso de la Probabilidad Clásica de Kolmogorov}
La toma de decisiones táctica y estratégica en entornos de mando y control (C2) se enfrenta de manera persistente a la denominada ``niebla de guerra'', caracterizada por información incompleta, ruidosa y manipulada por el adversario. Los marcos computacionales tradicionales de soporte a la decisión (POMDP, MILP y Redes Bayesianas) se fundamentan axiomáticamente en la probabilidad clásica de Kolmogorov.

Estructurar espacios de decisión complejos bajo este paradigma presenta limitaciones estructurales severas:
\begin{enumerate}
    \item \textbf{Asunción de Conmutatividad ($\text{Im}(z) = 0$):} Los modelos clásicos asumen que el orden de evaluación u observación de factores independientes no altera la distribución de probabilidad del estado final ($[\hat{A}, \hat{B}] = 0$). En la práctica operacional, las intervenciones tácticas no conmutan; la secuencia temporal de las acciones genera trayectorias de fase no equivalentes.
    \item \textbf{Incapacidad para Modelar Interferencia Contextual:} Tratan la ambigüedad como varianza estocástica en lugar de términos de fase capaces de generar interferencias constructivas o destructivas.
    \item \textbf{Colapso ante Marginales Indistinguibles:} Cuando las señales individuales recopiladas por sensores aislados parecen ruido blanco, los métodos estocásticos convencionales colapsan hacia la adivinación aleatoria.
\end{enumerate}

\section{El Problema del Aliasing Relacional (CSRA)}
Para acotar la clase de problemas donde la arquitectura clásica colapsa, definimos formalmente la clase de problemas \emph{Context-Sparse Relational Aliasing} (CSRA). Un escenario CSRA se caracteriza por:
\begin{itemize}
    \item Un conjunto de $P$ relaciones candidatas entre variables, de las cuales solo una es contextualmente relevante en cada evento.
    \item Distribuciones marginales de primer orden individualmente idénticas e indistinguibles entre las clases latentes.
    \item Información útil codificada exclusivamente en la orientación de una relación de segundo orden.
    \item Una señal de contexto que indica qué relación debe consultarse.
\end{itemize}

Bajo estas condiciones, los modelos genéricos de aprendizaje automático colapsan a un rendimiento cercano al azar ($\sim 27.26\%$ de precisión). En contraste, QAF extrae un estadístico relacional contextual $[c_R, c_I]$ equivalente a la parte real e imaginaria de la coherencia compleja de la 2-RDM, alcanzando una precisión media del \textbf{96.96\%}.

\section{La Hipótesis del Marco Relacional y el Sesgo Inductivo}
La ventaja de QAF no proviene de procesar información en hardware físico cuántico, sino de operar con un \textbf{sesgo inductivo (\emph{inductive bias}) y una representación estructural superior}.

Esta afirmación se sostiene mediante dos pruebas de control falsables:
\begin{enumerate}
    \item \textbf{Control del Oráculo Clásico (\emph{Oracle-Gated-Relational}):} Cuando un clasificador clásico recibe exactamente el mismo estadístico suficiente de coherencia relacional seleccionada por contexto, el modelo clásico empata con precisión idéntica ($96.96\%$) al motor QAF.
    \item \textbf{Control Negativo de Patrones Visibles:} En datos con clases separables en primer orden, los algoritmos clásicos como PCA-KMeans dominan alcanzando un índice ARI de $0.905$, mientras que QAF queda cerca de cero ($ARI \approx 0$). Esto demuestra que QAF es una herramienta especializada para regímenes relacionales no evidentes.
\end{enumerate}

\section{Genealogía Intelectual: Las 9 Tradiciones}
QAF sintetiza nueve tradiciones históricas clave en una arquitectura operacional unificada:
\begin{table}[h]
\centering
\caption{Genealogía Intelectual del Quantum Analysis Framework (QAF)}
\vspace{0.3cm}
\begin{tabular}{lll}
\toprule
\textbf{Tradición / Autor} & \textbf{Contribución Clave} & \textbf{Herencia en QAF} \\
\midrule
Newton (1687) & Campo y potencial ($F = G \frac{m_1 m_2}{r^2}$) & Campo de relevancia contextual \\
Gauss (1809) & Inferencia geométrica y error & Calibración de evidencia e incertidumbre \\
Ising (1925) & Redes de espines acoplados ($J_{ij}$) & Matriz de tensiones entre decisiones \\
Dantzig (1963) & Programación lineal y región factible & Optimización bajo restricciones de misión \\
Bellman (1957) & Programación dinámica multietapa & Continuidad temporal y memoria de estado \\
Hopfield (1982) & Atractores en paisajes de energía & Convergencia a estados de coherencia \\
Pearl (1988--1995) & Grafos causales e intervención $do(\cdot)$ & Verificación causal de acciones \\
Tradición QUBO & Unificación binaria ($\min_x x^T Q x$) & Lenguaje de compilación discreta \\
Tensor Networks & Compresión tensorial (MPS/QTT) & Motor de compresión y transporte de fase \\
\bottomrule
\end{tabular}
\end{table}


\chapter{Espacios de Decisión No Conmutativos y Teoría de Decisiones Cuánticas (QDT)}
\label{cap:qdt_no_conmutativo}

\section{Construcción del Espacio de Hilbert Operacional}
Para superar las restricciones conmutativas de la probabilidad clásica de Kolmogorov sin recurrir a embeddings no físicos, QAF mapea el espacio de decisión a un espacio de Hilbert vectorial $\mathcal{H}$.

\subsection{Mapeo Directo para Espacios Combinatorios Discretos}
Para un sistema de $K$ variables binarias $x_a \in \{0, 1\}$, QAF establece una isometría $1:1$ hacia el espacio de Hilbert local $\mathcal{H}_a \cong \mathbb{C}^2$:
\begin{equation}
\mathcal{H} = \bigotimes_{a=1}^K \mathcal{H}_a \cong (\mathbb{C}^2)^{\otimes K}
\end{equation}
Cada plan discreto $x = (x_1, \dots, x_K)^T$ se asocia al ket base $|x_1 x_2 \dots x_K\rangle$. El estado global $|\psi(t)\rangle$ evoluciona como una superposición continua de dichos planes:
\begin{equation}
|\psi(t)\rangle = \sum_{x_1, \dots, x_K \in \{0,1\}} c_{x_1 \dots x_K}(t) |x_1 x_2 \dots x_K\rangle
\end{equation}
donde $c_{x_1 \dots x_K}(t) \in \mathbb{C}$ representa la amplitud de probabilidad compleja y la fase contextual asociada al plan $x$.

\section{Álgebra de Operadores No Conmutativos}
Las decisiones se representan mediante operadores hermíticos construidos con álgebra de Pauli y proyectores locales:
\begin{equation}
\hat{n}_a = \frac{1}{2}(I - \sigma_z^{(a)})
\end{equation}
El conmutador entre dos decisiones $\hat{A}_a$ y $\hat{B}_b$ mide la sensibilidad a la secuencia temporal de intervenciones:
\begin{equation}
[\hat{A}_a, \hat{B}_b] = \hat{A}_a \hat{B}_b - \hat{B}_b \hat{A}_a \neq 0
\end{equation}
La no conmutatividad obliga al sistema a mantener trayectorias de fase no equivalentes según la secuencia de acciones, un fenómeno inaccesible para modelos Bayesianos o Markovianos clásicos.

\section{Formulación de Probabilidad Cuántica (QDT)}
Según la Teoría de Decisiones Cuánticas (QDT), la probabilidad $p(\pi_n, t)$ de seleccionar una alternativa $\pi_n$ se compone de dos términos:
\begin{equation}
p(\pi_n, t) = f(\pi_n, t) + q(\pi_n, t)
\end{equation}
donde $f(\pi_n, t)$ es el factor de utilidad objetiva clásica:
\begin{equation}
f(\pi_n, t) = \prod_{a=1}^K \text{Tr}\left( \rho^{(a)}(t) \hat{n}_a^{x_a} (1 - \hat{n}_a)^{1-x_a} \right)
\end{equation}
y $q(\pi_n, t) \in [-1, 1]$ es el factor de atracción e interferencia contextual derivado de la 2-RDM:
\begin{equation}
q(\pi_n, t) = 2 \sum_{a<b} \text{Re}\left( \rho^{(a,b)}_{x_a x_b, \bar{x}_a \bar{x}_b}(t) \right)
\end{equation}

\section{La Fase Compleja y la Analogía de la Impedancia Eléctrica}
El \emph{inductive bias} del QAF se comprende mediante la analogía de la impedancia en circuitos de corriente alterna ($Z = R + jX$):
\begin{itemize}
    \item \textbf{Componente Real / Resistiva ($c_R = \text{Re}(z)$):} Mide la coincidencia directa de amplitud y alineación relacional. La IA clásica opera exclusivamente en este dominio real.
    \item \textbf{Reactancia Inductiva Positiva ($c_I > 0 \leftrightarrow +jX_L$ / Inductor):} Como un inductor que almacena energía en el campo magnético y avanza la fase, $c_I > 0$ representa la \textbf{interferencia constructiva} y atracción táctica hacia relaciones auténticas.
    \item \textbf{Reactancia Capacitiva Negativa ($c_I < 0 \leftrightarrow -jX_C$ / Condensador):} Como un condensador que almacena energía en el campo eléctrico y retarda la fase, $c_I < 0$ representa la \textbf{interferencia destructiva} y repulsión táctica ante engaños y señuelos (\emph{spoofing}).
\end{itemize}
Al incorporar $c_I$ en su sesgo inductivo, QAF detecta cuando un señuelo o interferencia genera una fase destructiva ($c_I < 0$), eliminando falsas alarmas que engañan a los modelos puramente resistivos.


% ==========================================
% PARTE II
% ==========================================
\part{El Motor Tensorial (MPS / QTT) y la Dinámica Variacional}

\chapter{Compresión del Espacio de Estados mediante Redes Tensoriales (MPS / QTT) y Reordenamiento Topológico (RCM)}
\label{cap:redes_tensoriales_rcm}

\section{Bypasando la Explosión Exponencial $2^K$: Representación MPS / QTT}
Para gestionar un espacio de decisión parametrizado por $K$ variables binarias discretas $x_a \in \{0, 1\}$, la representación densa convencional exige almacenar un vector de estados de dimensión $2^K$. Para $K = 500$ variables, un vector denso requeriría $\sim 10^{146}$ GB de memoria RAM, haciendo inviable la evaluación en tiempo real.

El Quantum Analysis Framework (QAF) supera esta barrera descomponiendo el vector de estado global $|\psi(t)\rangle$ en una cadena unidimensional (1D) de $K$ tensores mediante Matrix Product States (MPS) o Quantized Tensor Trains (QTT):
\begin{equation}
|\psi(t)\rangle = \sum_{x_1, \dots, x_K \in \{0,1\}} A^{x_1} A^{x_2} \cdots A^{[K]x_K} |x_1 x_2 \dots x_K\rangle
\end{equation}
donde cada $A^{[a]x_a}$ es un tensor local de rango 3 asociado al sitio ($site$) $a$, con dimensión física $d = 2$ e índices de enlace virtuales $\alpha_{a-1}, \alpha_a$ de dimensión acotada $\chi$ (\emph{bond dimension}).

\section{Cota de Memoria y Escalabilidad de la Dimensión de Enlace ($\chi$)}
La dimensión de enlace virtual $\chi$ mide la capacidad de la red para almacenar entrelazamiento y correlaciones de fase. Acotando $\chi \le 64$, el consumo de memoria del motor QAF escala de forma estrictamente lineal/polinómica:
\begin{equation}
\text{Memoria MPS} = \mathcal{O}(K \cdot d \cdot \chi^2)
\end{equation}

\subsection{Impacto Cuantitativo de la Compresión}
\begin{itemize}
    \item \textbf{Para $K = 512$ variables ($\chi = 16$):} La memoria pasa de $2^{512}$ bytes a solo \textbf{0.56 MB}, permitiendo ejecutar el motor en procesadores embebidos de a bordo.
    \item \textbf{Para $K = 50.000.000$ variables ($\chi = 16$):} El estado tensorial se almacena en \textbf{204.8 GB} de RAM, frente a los 20 Petabytes que requeriría una matriz densa.
\end{itemize}

\section{Reordenamiento Topológico mediante Reverse Cuthill-McKee (RCM)}
Dado que la topología física de una red MPS es 1D, si dos variables altamente acopladas quedasen separadas por muchos sitios intermedios, la dimensión de enlace $\chi$ crecería desproporcionadamente para transportar el entrelazamiento.

Para prevenir la saturación de $\chi$, QAF ejecuta un reordenamiento topológico previo con el algoritmo Reverse Cuthill-McKee (RCM):
\begin{enumerate}
    \item Construye el grafo de interacciones $G(V, E)$ a partir de la matriz de costos.
    \item Minimiza el ancho de banda $\beta(G) = \max_{(a,b) \in E} |\pi(a) - \pi(b)| \le b \le 8$.
    \item Reordena los tensores en la cadena 1D para que las variables acopladas sean vecinas físicas inmediatas.
\end{enumerate}
Esto garantiza que $\chi \le 64$ permanezca acotada y genera matrices QUBO bandejadas que se compilan en solo \textbf{18.5 ms} para $K = 500$ variables.


\chapter{Evolución Variacional Continua (1TDVP y 2TDVP Adaptativo ante Shocks)}
\label{cap:tdvp_dinamica}

\section{El Hamiltoniano Operacional en Representación MPO}
La dinámica temporal del espacio de decisión está gobernada por un Hamiltoniano operacional representado como un Matrix Product Operator (MPO). La estructura del MPO codifica de forma compacta tanto los costos individuales como las interacciones complejas:
\begin{itemize}
    \item \textbf{Términos Diagonales ($H_{aa}$):} Codifican el costo individual de activar la variable $a$ (prioridad táctica o consumo de recursos).
    \item \textbf{Términos Fuera de Diagonal ($H_{ab}$):} Codifican acoplamientos cuadráticos, incompatibilidades físicas y penalizaciones por saturación.
\end{itemize}

\section{Integración Variacional con 1TDVP (Projector-Splitting)}
Para evolucionar $|\psi(t)\rangle$ sin abandonar la variedad MPS de \emph{bond dimension} $\chi$ acotada, QAF utiliza el algoritmo 1-site TDVP mediante \emph{projector-splitting}:
\begin{equation}
i\hbar_{\text{eff}} \frac{d}{dt} |\psi(t)\rangle = P_{T_{|\psi\rangle} \mathcal{M}_\chi} H(t) |\psi(t)\rangle
\end{equation}
\begin{enumerate}
    \item \textbf{Conservación Unitaria:} Preserva la norma ($\langle\psi(t)|\psi(t)\rangle = 1$) y la estructura de calibre de la red.
    \item \textbf{Escalabilidad:} Complejidad por barrido de $\mathcal{O}(K \cdot d^3 \cdot \chi^3)$, estrictamente lineal con respecto a $K$.
\end{enumerate}

\section{Dinámica Adaptativa 2TDVP ante Shocks de Entrelazamiento}
Ante incrementos drásticos de entrelazamiento ($d\chi/dt > \epsilon$), QAF conmuta dinámicamente al algoritmo 2-site TDVP (2TDVP):
\begin{equation}
\chi_{\text{adapt}}(t) = \min\left(\chi_{\text{max}}, \text{rank}_\eta(\Theta)\right)
\end{equation}
Si el entrelazamiento aumenta, $\chi$ se expande dinámicamente hasta $\chi_{\text{max}} = 64$; cuando el sistema se estabiliza, $\chi$ se contrae automáticamente para ahorrar memoria, manteniendo una fidelidad exacta de fase ($1.000000$).


% ==========================================
% PARTE III
% ==========================================
\part{Extracción de Fase y Compilación Discreta}

\chapter{Operadores de Densidad Reducida (1-RDM / 2-RDM) y Plano de Fase Complejo}
\label{cap:rdm_plano_fase}

\section{Extracción de Densidad Reducida Mediante Contracción de Entornos MPS}
En QAF, la evaluación de decisiones y la extracción de correlaciones no requiere reconstruir el vector global de dimensión $2^K$, sino realizar contracciones locales con entornos $L^{[a]}$ y $R^{[a+1]}$:
\begin{enumerate}
    \item \textbf{1-RDM ($\rho^{(a)}$ - $2 \times 2$):} Trazo de $K-1$ sitios para obtener la probabilidad marginal de activación $\langle n_a \rangle = \text{Tr}(\rho^{(a)} \hat{n}_a)$.
    \item \textbf{2-RDM ($\rho^{(a,b)}$ - $4 \times 4$):} Contracción del tensor fusionado $\Theta_{a,b}$ para aislar la coherencia relacional fuera de la diagonal ($\rho^{(a,b)}_{00,11}$).
\end{enumerate}

\section{El Cociente de Coherencia Relacional Compleja $z_{ab} = c_R + i c_I$}
A partir de la 2-RDM, QAF define el parámetro de coherencia relacional en el círculo unitario ($|z_{ab}| \le 1$):
\begin{equation}
z_{ab} = \frac{\rho^{(a,b)}_{00,11}}{\sqrt{\rho^{(a,b)}_{00,00} \cdot \rho^{(a,b)}_{11,11}}} = c_R + i c_I
\end{equation}

\section{Interpretación en el Plano de Fase: La Analogía del Circuito AC}
El parámetro $z_{ab} = c_R + i c_I$ equivale a la impedancia $Z = R + jX$:
\begin{itemize}
    \item \textbf{Parte Real / Resistiva ($c_R = \text{Re}(z)$):} Coincidencia directa de amplitud. Dominio único de la IA clásica ($X = 0$).
    \item \textbf{Reactancia Inductiva ($c_I > 0 \leftrightarrow +jX_L$ / Inductor):} Interferencia constructiva y atracción táctica hacia relaciones auténticas.
    \item \textbf{Reactancia Capacitiva ($c_I < 0 \leftrightarrow -jX_C$ / Condensador):} Interferencia destructiva y repulsión táctica ante engaños y señuelos.
\end{itemize}

\section{Mitigación de Señuelos y Aislamiento de Falsas Alarmas (Spoofing)}
Ante señuelos adversarios de igual amplitud individual, QAF detecta la reactancia capacitiva destructiva ($c_I < 0$), neutralizando el señuelo y reduciendo el desperdicio de recursos defensivos en un \textbf{53,91\%}.


\chapter{El Compilador QUBO Bandejado y Ejecución Discreta}
\label{cap:compilador_qubo}

\section{Muestrono Asíncrono Disparado por Eventos}
El estado tensorial $|\psi(t)\rangle$ se congela y proyecta hacia la compilación discreta únicamente bajo tres condiciones gatillo:
\begin{enumerate}
    \item \textbf{Shock de Gradiente:} $\|dH/dt\| > \epsilon$ (aparición de amenazas o pérdida de nodos).
    \item \textbf{Estabilización de Información:} $dI/dt \to 0$ (fase variacional coherente).
    \item \textbf{Temporizador Táctico:} $\Delta t > \Delta t_s$ (frecuencia mínima de refresco).
\end{enumerate}

\section{Derivación Analítica de la Matriz QUBO $Q(t_k)$}
En el instante $t_k$, el compilador traduce $|\psi(t_k)\rangle$ a la función de energía binaria $E(x) = x^T Q(t_k) x$:
\begin{align}
Q_{aa} &= H_{aa} \cdot \langle n_a \rangle \\
Q_{ab} &= H_{ab} \cdot \text{Re}(z_{ab}) \cdot \sqrt{\langle n_a \rangle \langle n_b \rangle}
\end{align}
Cuando $c_I = \text{Im}(z_{ab}) < 0$ (fase destructiva), el compilador penaliza automáticamente $Q_{ab}$, eliminando asignaciones a señuelos.

\section{Estructura Bandejada y Mitigación de Minor-Embedding}
El reordenamiento RCM previo garantiza una matriz QUBO bandejada ($b \le 8$):
\begin{itemize}
    \item \textbf{Reducción de Memoria:} Para $K = 50.000.000$ variables, la memoria QUBO se reduce de 20 Petabytes a solo \textbf{14 GB} ($\mathcal{O}(K \cdot b)$).
    \item \textbf{Mitigación de Minor-Embedding:} Elimina cadenas largas de qubits en \emph{annealers} cuánticos (D-Wave), previniendo el colapso por ruido térmico.
    \item \textbf{Ejecución en Tiempo Real:} Permite a los \emph{solvers} discretos resolver $K = 500$ variables en \textbf{18.5 ms}.
\end{itemize}


% ==========================================
% PARTE IV
% ==========================================
\part{Inteligencia Artificial Metacognitiva y Multiagente}

\chapter{El Enrutador Metacognitivo QAFRouter e Inteligencia Artificial Multiagente}
\label{cap:qaf_router}

\section{Racionalidad de Recursos en Sistemas Multiagente}
El patrón de comunicación global \emph{Broadcast-All} resulta insostenible en entornos de mando y control (C2), pues genera consumo redundante de tokens, saturación del ancho de banda y amplificación del sesgo de mayoría correlacionada.

QAF integra la capa metacognitiva \texttt{QAFRouter}, basada en la racionalidad de recursos (\emph{resource rationality}), que evalúa el estado de misión y activa solo a los agentes especializados necesarios con el contexto mínimo suficiente.

\section{Evaluación de las Variables de Estado de Misión}
El \texttt{QAFRouter} evalúa continuamente cuatro métricas de estado derivadas del motor MPS/2-RDM:
\begin{equation}
G_i = \omega_1 U + \omega_2 C + \omega_3 R + \omega_4 G
\end{equation}
donde $U$ es Incertidumbre (entropía), $C$ es Contradicción (divergencia entre sensores), $R$ es Riesgo Operacional e $G$ es Brecha de Información.

\section{Funciones de Activación Condicional y Puertas de Seguridad}
\begin{itemize}
    \item \textbf{Memoria y Recuperación:} Activado con $U$ moderada y $G$ elevada.
    \item \textbf{Especialista / Verificador:} Activado ante alta contradicción $C$, aplicando el filtro de fase 2-RDM para aislar señuelos (CSRA).
    \item \textbf{Puerta Dura de Seguridad (\emph{Hard Safety Gate}):} Si $R \ge 0.82$, bloquea la delegación autónoma y exige supervisión humana directa.
\end{itemize}

\section{Impacto en Eficiencia: Reducción de Tokens y Tráfico}
\begin{itemize}
    \item \textbf{Reducción del 62.8\% en Tokens Totales:} Caída del $64.7\%$ en contexto de entrada frente a \emph{Broadcast-All}.
    \item \textbf{Reducción del 69.87\% en Tráfico de Mensajes:} Optimización masiva para redes de radio tácticas degradadas.
    \item \textbf{Mejora del 17.88\% en Éxito de Coordinación:} Bajo escenarios con hasta un $60\%$ de interferencia o ruido adversario.
\end{itemize}


\chapter{Integración Práctica: Control de Enjambres de Drones Autónomos y C2}
\label{cap:enjambres_drones_c2}

\section{Percepción de IA y Generación de Priors Semánticos a $t=0$}
El motor tensorial QAF opera acoplado con la capa de percepción basada en Inteligencia Artificial y visión por computador. En el instante $t=0$, los modelos de visión procesan sensores EO/IR y radar para generar priors semánticos que alimentan los términos diagonales del MPO ($H_{aa}$).

\section{Filtrado de Fase 2-RDM en Entornos de Guerra Electrónica / Spoofing}
Ante señuelos electromagnéticos que imitan la amplitud de amenazas reales, QAF aplica el filtro de fase relacional de la 2-RDM:
\begin{itemize}
    \item \textbf{Reactancia Capacitiva ($c_I < 0$):} Detecta la descorrelación de fase del señuelo con la red de sensores.
    \item \textbf{Anulación en QUBO:} Cancela los términos cruzados $Q_{ab}$, reduciendo el desperdicio de interceptores en señuelos en un \textbf{53.91\%}.
\end{itemize}

\section{Resiliencia y Reconfiguración Autónoma ante Pérdida de Nodos (50\%)}
Ante la destrucción imprevista del $50\%$ del enjambre:
\begin{enumerate}
    \item \textbf{Shock de Gradiente:} Congela el MPS al cumplir $\|dH/dt\| > \epsilon$.
    \item \textbf{RCM Dinámico:} Recalcula $\pi^*$ sobre el grafo superviviente, manteniendo $b \le 8$.
    \item \textbf{Re-compilación Instantánea:} Genera la nueva asignación en \textbf{18.5 ms} sin colapso centralizado.
\end{enumerate}


% ==========================================
% PARTE V
% ==========================================
\part{Demostraciones Empíricas y Escalabilidad a 50 Millones}

\chapter{El Benchmark CSRA y la Ventaja Estructural (96.96\% vs 27.26\%)}
\label{cap:benchmark_csra}

\section{Definición Formal de la Clase de Problemas CSRA}
Para acotar el régimen de datos donde los modelos clásicos colapsan, se define la clase de problemas \emph{Context-Sparse Relational Aliasing} (CSRA) bajo cuatro condiciones:
\begin{enumerate}
    \item \textbf{Solapamiento Marginal:} Distribuciones de primer orden idénticas e indistinguibles entre clases (ruido blanco).
    \item \textbf{Información en Fase de 2º Orden:} Señal útil codificada en la orientación relacional entre pares.
    \item \textbf{Señuelos Masivos:} Pares irrelevantes con igual magnitud y energía.
    \item \textbf{Contexto Disperso:} Señal indicadora de la relación relevante a consultar.
\end{enumerate}

\section{Resultados del Benchmark CSRA y Ventaja Estructural}
Los clasificadores tradicionales (RBF-SVM, ExtraTrees, Random Forest) colapsan cerca del azar ($\sim 27.26\%$ de precisión). En contraste, QAF extrae la coherencia relacional $z_{ab} = c_R + i c_I$ de la 2-RDM, alcanzando una precisión media del \textbf{96.96\%} con solo 2 variables.

\section{Prueba de Ablación y Demostración Inductiva}
La prueba con el oráculo \texttt{Oracle-Gated-Relational} demuestra que cuando un modelo clásico recibe el estadístico $[c_R, c_I]$, alcanza el mismo \textbf{96.96\%}. Esto prueba que la ganancia es $100\%$ representacional, derivada de un sesgo inductivo superior para medir fases complejas.


\chapter{Asignación Táctica C2 y Mitigación de Señuelos}
\label{cap:asignacion_c2_senuelos}

\section{Escenario Táctico de Defensa Aérea ($K = N \times M$)}
En sistemas de Mando y Control (C2), la asignación de $N$ interceptores a $M$ amenazas genera un espacio combinatorio $K = N \times M$ caracterizado por la presencia de señuelos descorrelacionados (\emph{spoofed decoys}).

\section{Reducción del Desperdicio de Recursos en Señuelos (53.91\%)}
Los algoritmos clásicos (MILP, Greedy) evalúan magnitudes aisladas y malgastan interceptores en señuelos. QAF inspecciona la reactancia capacitiva de fase en la 2-RDM ($c_I = \text{Im}(z_{ab}) < 0$), reduciendo el desperdicio de recursos defensivos del $13.26\%$ al $6.11\%$ (\textbf{reducción relativa del 53.91\%}).

\section{Calibración de Confianza y Validación con Oráculo ($K = 8$)}
\begin{enumerate}
    \item \textbf{Mejora en Brier Score:} Reducción de $0.1243$ a \textbf{0.1079} ($13.22\%$ de mejora en calibración de incertidumbre).
    \item \textbf{Fidelidad Exacta a $K = 8$:} Fidelidad de \textbf{1.000000} frente al oráculo denso de Schrödinger, confirmando la precisión de la compresión MPS/TDVP.
\end{enumerate}


\chapter{Escalabilidad Estructural de K=8 a 50 Millones de Variables}
\label{cap:escalabilidad_50m}

\section{Análisis Comparativo de Regímenes de Memoria}
El escalado computacional de QAF desde $K = 8$ hasta $K = 50.000.000$ demuestra la viabilidad práctica de la compresión tensorial:
\begin{itemize}
    \item \textbf{Régimen Denso Clásico ($2^K$ / $K^2$):} Requeriría $\sim 20\text{ Petabytes}$ de RAM para 50M variables, resultando inviable.
    \item \textbf{Régimen Disperso Bandejado CSR ($\mathcal{O}(K \cdot b)$):} Gracias al reordenamiento RCM ($b \le 8$), la matriz QUBO requiere solo \textbf{14 GB} de RAM.
    \item \textbf{Régimen Tensorial MPS/QTT ($\mathcal{O}(K \cdot d \cdot \chi^2)$):} Con $\chi = 16$, almacena el estado de 50M variables en \textbf{204.8 GB} de RAM (y solo \textbf{0.56 MB} para $K=500$).
\end{itemize}

\section{Tiempos de Ejecución y Amortización por Variable}
El motor variacional 1TDVP/2TDVP presenta complejidad temporal estrictamente lineal por barrido: $\mathcal{O}(K \cdot d^3 \cdot \chi^3)$.
Para $K = 500$ variables, la compilación QUBO se ejecuta en \textbf{18.5 ms}, con un costo amortizado de $\sim 23.5\text{ ns}$ por variable.

\section{Prueba Streaming por Bloques a 50 Millones de Variables}
En pruebas empíricas de compilación en \emph{streaming} a 50M de variables:
\begin{itemize}
    \item \textbf{Tiempo Total:} \textbf{1.07 segundos}.
    \item \textbf{Rendimiento (\emph{Throughput}):} \textbf{45.5 millones de variables/s}.
    \item \textbf{Ajuste Lineal:} Coeficiente $R^2 = 0.998$, confirmando la eliminación de cuellos de botella exponenciales.
\end{itemize}


% ==========================================
% PARTE VI
% ==========================================
\part{Guía del Desarrollador y Repositorio de Código (GitHub)}

\chapter{Arquitectura del Repositorio Abierto y Guía del Desarrollador en GitHub}
\label{cap:github_desarrollador}

\section{Estructura Modular del Repositorio (\texttt{qaf-engine})}
Para garantizar la máxima reproducibilidad científica, el código fuente del Quantum Analysis Framework (QAF) se organiza en el repositorio modular \texttt{qaf-engine}:

\begin{verbatim}
qaf-engine/
|-- README.md                      # Visión general, instalación y guía rápida
|-- LICENSE                        # Licencia dual (MIT / Apache 2.0)
|-- docs/                          # Documentación LaTeX del libro y manuales
|-- src/                           # Código fuente del motor
|   |-- core/                      # Redes tensoriales: MPS, MPO, 1TDVP, 2TDVP
|   |-- rdm/                       # Extracción 1-RDM / 2-RDM y fase z
|   |-- compiler/                  # Reordenamiento RCM y compilador QUBO
|   \-- router/                    # QAFRouter metacognitivo para multiagente
|-- benchmarks/                    # Scripts de validación reproducible
|   |-- csra_benchmark.py          # Benchmark de ventaja relacional CSRA
|   |-- c2_allocation.py           # Asignación táctica y mitigación de señuelos
|   |-- token_efficiency.py        # Evaluación de footprint de tokens
|   \-- scaling_50m.py             # Benchmark streaming de 1M a 50M
|-- tests/                         # Pruebas unitarias y oráculos (K=8)
\-- notebooks/                     # Jupyter Notebooks demostrativos e interactivos
\end{verbatim}

\section{Batería de Pruebas Unitarias e Integración Continua (CI/CD)}
El repositorio incorpora integración continua mediante GitHub Actions:
\begin{enumerate}
    \item \textbf{Verificación con Oráculo ($K=8$):} Exige una fidelidad superior a $0.999999$ frente a la simulación densa exacta.
    \item \textbf{Tests de Reordenamiento RCM:} Verifica la reducción del ancho de banda a $b \le 8$.
    \item \textbf{Conservación Unitaria en 1TDVP:} Garantiza $|\langle\psi|\psi\rangle - 1| < 10^{-12}$ a lo largo de toda la integración.
\end{enumerate}

\section{Hoja de Ruta de Desarrollo Futuro (Fase II)}
\begin{itemize}
    \item \textbf{Producción en Julia (\texttt{ITensors.jl}):} Optimización del núcleo tensorial para reducir latencias JIT.
    \item \textbf{Aceleración GPU con NVIDIA cuQuantum:} Integración de \texttt{cuTensorNet} para clústeres A100/H100.
    \item \textbf{Ganancia Adaptativa por Entropía $\gamma(U)$:} Ajuste dinámico de fase bajo alta incertidumbre.
\end{itemize}

\end{document}
